\section{Four-dimensional inducing blocks and their remaining fields}
\label{sec:rank-four-inducing}

\begin{theorem}[D-BI-11: a rank-four field obstruction]
\label{D-BI-11}
Let a finite-cover Weil sixfold source have exact imaginary quadratic
character field $K$, period rank nine, and nonsplit actual polarization.
Suppose its complex coefficient representation is induced from an
irreducible representation $\tau$ of dimension four.  Then
\[
 K\in\{\mathbf Q(\sqrt{-d}):d=1,2,3,5,6,15\}.
                                                               \tag{BI.23}
\]
In particular no four-dimensional inducing block can occur for
\[
                    d=7,11,19,43,67,163.                       \tag{BI.24}
\]
This is an all-degree assertion about the original representation;
the inducing character field is not assumed to be quadratic.

Among the eight non-Eisenstein class-number-one fields, the only
remaining possibilities have $K=\mathbf Q(i)$ or
$\mathbf Q(\sqrt{-2})$.  Write $L=\mathbf Q(\chi_\tau)$.
In these remaining cases $L$ is a proper biquadratic extension of
$\mathbf Q$, $[L:K]=2$, and its possible values are
\[
\begin{array}{c|c}
 K& L\\\hline
 \mathbf Q(i)&\mathbf Q(\zeta_8),\ \mathbf Q(\zeta_{12}),\
                         \mathbf Q(i,\sqrt5)\\
 \mathbf Q(\sqrt{-2})&\mathbf Q(\zeta_8).
\end{array}                                                    \tag{BI.25}
\]
Every such remaining case has exactly twelve branch points on the
intermediate rational base: ten reflection inertias and two inertias
of moving rank two outside the reflection normal closure.
Their quotient by that normal closure is cyclic and its covering
curve has genus zero.  These are necessary conditions, not a claim
that one of the remaining families exists or has rank nine.
\end{theorem}
\begin{proof}
Write the original complex representation as
$\rho=\operatorname{Ind}_H^G\tau$ and let $J$ be the faithful
image of $\tau$.  If $\tau$ is imprimitive, a system of blocks in
dimension four has block dimension one or two.  Transitivity of
induction then makes $\rho$ respectively monomial, or induced
from a two-dimensional representation.  D-MON-01--03 and D-BI-01
exclude these cases.  We may assume $\tau$ primitive.

Pass to the actual connected $J$-cover $C_0\to X_0$ obtained by
quotienting the kernel in the intermediate $H$-cover.  Frobenius
reciprocity identifies the complex multiplicity variation with
the original one, including its period differential kernel, as
in D-BI-01.  If $h=g(X_0)$ and $D$ is the total moving dimension,
then $D=14-8h$.  At $h=1$ the period variation factors through
at most six marked-base parameters, and $h\ge2$ is impossible.
Thus rank nine forces $h=0$ and
\[
 r\ge12,\qquad D=14,\qquad
 \sum_{j=1}^r\bigl(\operatorname{rank}(1-g_j)-1\bigr)=14-r\le2.
                                                               \tag{BI.26}
\]
In particular at least ten of the generating inertia matrices are
reflections, and no inertia has moving rank four.

Let $R\lhd J$ be the subgroup generated by all reflections of
$J$.  It is nontrivial.  Its fixed subspace is $J$-stable and
therefore zero.  The nonorthogonality components of its reflection
root lines give an orthogonal decomposition permuted by $J$.
Primitivity implies that there is only one component, so $R$ is
an irreducible rank-four reflection group.

There are at most two inertia matrices outside $R$.  There cannot
be exactly one, since their product in $J/R$ is one.  If there are
none, the inertia generators give $J=R$.  Otherwise there are
exactly two; (BI.26) forces their moving ranks to be two, $r=12$,
and every other inertia to be a reflection.  The quotient $J/R$
is generated by the coset of either outside inertia, and the other
coset is its inverse.  Hence it is cyclic and
\[
                         J=\langle R,g\rangle               \tag{BI.27}
\]
for either outside inertia $g$.  The quotient cover
$C_0/R\to\mathbf P^1$ has exactly two branch points, both with
full cyclic inertia, and Riemann--Hurwitz gives genus zero.

We now use the published rank-four reflection classification
\cite[Theorem~5.12, p.~435]{FIELD:Cohen}.  An imprimitive $R$ has
monomial normalizer except for $G(2,2,4)=W(D_4)$, by
\cite[Proposition~2.13, pp.~391--392]{FIELD:Cohen}.  The monomial
case is already excluded.  The primitive real possibilities are
$A_4,F_4,H_4$, and the nonreal primitive possibilities are
$G_{29},G_{31},G_{32}$.  The $A_4$ normalizer is the scalar
standard $\mathfrak S_5$ group, and is excluded by D-BI-08.
The remaining normalizer models can be chosen as follows:
\[
\begin{array}{c|c|c}
 R&\text{finite normalizer model }N\text{ up to scalars}
                                      &\text{matrix field }F\\\hline
 W(D_4)&W(F_4)&\mathbf Q\\
 W(F_4)&\langle W(F_4),\Phi/\sqrt2\rangle&\mathbf Q(\sqrt2)\\
 W(H_4)&W(H_4)&\mathbf Q(\sqrt5)\\
 G_{29},G_{31}&R&\mathbf Q(i)\\
 G_{32}&R&\mathbf Q(\sqrt{-3}).
\end{array}                                                    \tag{BI.28}
\]

Here the $D_4$ assertion is exactly Cohen's Proposition~2.13.
For $F_4$ the rational simple-root matrices and $\Phi$ are those
of D-GF4-01.  The matrix $\Phi/\sqrt2$ interchanges the two
root-length classes, squares to $-I$, and normalizes $W(F_4)$.
An orthogonal normalizer preserving the two classes is a root-system
automorphism and belongs to $W(F_4)$, by the chamber argument and
the absence of a nontrivial oriented-diagram automorphism.  This
gives the displayed finite group of order $2304$ and its field.
For $H_4$ all roots have one length and its Coxeter diagram has
no nonidentity automorphism, so the same chamber argument gives
$N=W(H_4)$.  Its standard roots and reflection matrices are over
$\mathbf Q(\sqrt5)$.  These real-normalizer reductions also follow
from Cohen's real-normalizer Lemma~2.12, p.~391, followed by the
root-system argument.  For the three nonreal groups the scalar
normalizer assertion is the \emph{corrected} Proposition~5.14:
\cite[p.~435]{FIELD:Cohen}\cite[p.~613]{FIELD:CohenErr}.
The original extra rank-four exception was deleted by that erratum.
The matrix fields are supplied by the standard reflection models
and \cite[Table~8.2, pp.~675--676]{FIELD:Malle}; only their role
as fields of definition is needed at this stage.

If $J=R$, its character field is contained in the displayed field
of its reflection model.  This either gives $K=\mathbf Q(i)$ or
$\mathbf Q(\sqrt{-3})$, or is incompatible with $K$ being imaginary.
In the remaining case write the outside inertia from (BI.27) as
\[
                           g=\lambda w,\qquad w\in N.
\]
Both $g$ and $w$ have finite order, so the commuting scalar
$\lambda$ is a root of unity.  Since $g$ fixes a two-dimensional
subspace, $\lambda^{-1}$ occurs with multiplicity two as an
eigenvalue of the matrix $w$ over $F$.  Its irreducible minimal
polynomial over $F$ consequently occurs at least twice in the
degree-four characteristic polynomial of $w$.  Therefore
\[
                 [F(\lambda):F]\le2,
                 \qquad K\subset L\subset F(\lambda).
                                                               \tag{BI.29}
\]
The first field inclusion is the induced-character identity; the
second follows from (BI.27) and the matrices in (BI.28).
There is no equality assumption between $K$, $L$, and $F(\lambda)$.

Here is the entire cyclotomic calculation in (BI.29).  If $F$ is
quadratic, the order $d$ of $\lambda$ has $\varphi(d)\le4$,
so $d\in\{1,2,3,4,5,6,8,10,12\}$.  This finite list follows
directly from $p-1\mid\varphi(d)$ for each prime $p\mid d$
and the prime-power formula for $\varphi$.  The degree condition
in (BI.29) gives exactly the following upper possibilities:
\[
\begin{array}{c|c|c}
 F&d&\text{imaginary quadratic subfields of all }F(\lambda)\\\hline
 \mathbf Q&1,2,3,4,6&-1,-3\\
 \mathbf Q(\sqrt2)&1,2,3,4,6,8&-1,-2,-3,-6\\
 \mathbf Q(\sqrt5)&1,2,3,4,5,6,10&-1,-3,-5,-15\\
 \mathbf Q(i)&1,2,3,4,6,8,12&-1,-2,-3\\
 \mathbf Q(\sqrt{-3})&1,2,3,4,6,12&-1,-3.
\end{array}                                                    \tag{BI.30}
\]
The entries in the last column denote the square-free radicands.
For example $\mathbf Q(\sqrt2,i)=\mathbf Q(\zeta_8)$ has
quadratic subfields with radicands $2,-1,-2$, whereas
$\mathbf Q(\sqrt2,\sqrt{-3})$ has radicands $2,-3,-6$.
For $F=\mathbf Q(\sqrt5)$ and $d=5,10$, the field is
$\mathbf Q(\zeta_5)$, whose sole quadratic subfield is the real
field $\mathbf Q(\sqrt5)$.  Every other degree-four field in
the table is biquadratic; its three quadratic subfields follow by
multiplying the two generating square classes.  This proves
(BI.23)--(BI.24).

We finish by establishing the proper-field assertion, which is
relevant to the actual polarization.  In this paragraph restrict
to $K=\mathbf Q(i)$ or $\mathbf Q(\sqrt{-2})$, and suppose $L=K$.
Frobenius reciprocity and irreducibility of $\rho$ give
\[
           \langle\tau,\operatorname{Res}_H^G\rho\rangle=1.
\]
The central $\tau$-projector of $K[H]$, acting on the given
$K$-model of $\rho$, therefore has a four-dimensional $K$-image
whose complexification is $\tau$.  Thus $\tau$ has a $K$-model;
no Schur-index assumption was silently made.  Its induced $K$-model
is isomorphic to the original one, and polarized Shapiro identifies
the actual multiplicity forms up to a positive rational scalar,
as in D-FIELD-03.  A scalar of this kind preserves splitness.
For $K=\mathbf Q(\sqrt{-2})$, the reflection-source theorem
D-FIELD-03 forces split determinant.  For $K=\mathbf Q(i)$,
the rank-four cases D-GAUSS-01 and D-G31-02 force the same
conclusion; the only Gaussian reflection component left by
D-GAUSS-03 has rank eight, not four.  This contradicts the
assumed nonsplit actual polarization.  Therefore $L\ne K$.

For $K=\mathbf Q(i)$ or $\mathbf Q(\sqrt{-2})$, inspection of
(BI.30) now gives precisely (BI.25).  Each ambient field containing
one of these $K$ is either that quadratic field itself or one of
the listed degree-four fields; once $K\subsetneq L$, no intermediate
choice remains.  This proves the theorem, including the exact
twelve-branch profile for all remaining cases.
\end{proof}
